\documentclass[10pt,a4paper]{article}

\usepackage[T1]{fontenc}
\usepackage[utf8]{inputenc}
\usepackage{lmodern}
\usepackage{amsmath,amssymb,amsthm}
\usepackage{booktabs,array}
\usepackage{geometry}
\usepackage{xcolor}
\usepackage{enumitem}
\usepackage{microtype}
\usepackage{hyperref}
\usepackage{fancyhdr}

\geometry{margin=1in}
\hypersetup{
  colorlinks=true,
  linkcolor=blue!55!black,
  citecolor=blue!55!black,
  urlcolor=blue!55!black,
  pdftitle={The Writing Was on the Floor: A Zero-Query Fresh-Message Forgery against GLYPH},
  pdfauthor={Martin Feussner}
}
\setlist{nosep,leftmargin=*}
\setlength{\parindent}{1em}
\setlength{\parskip}{0.25em}
\pagestyle{fancy}
\fancyhf{}
\fancyhead[L]{The Writing Was on the Floor}
\fancyhead[R]{13 October 2026}
\fancyfoot[C]{\thepage}
\renewcommand{\headrulewidth}{0.3pt}

\newtheorem{theorem}{Theorem}
\newtheorem{lemma}{Lemma}
\newcommand{\R}{\mathcal{R}}
\newcommand{\Rq}{\mathcal{R}_q}
\newcommand{\Kfloor}[1]{\left\lfloor #1\right\rfloor_K}
\newcommand{\Hash}{\mathsf{H}}
\newcommand{\Encode}{\mathsf{F}}

\title{\bfseries The Writing Was on the Floor: A Zero-Query Fresh-Message Forgery against GLYPH}
\author{Martin Feussner\\
\small Selmer Center, University of Bergen\\
\small \texttt{martin.feussner@uib.no}}
\date{13 October 2026}

\begin{document}
\maketitle

\noindent\fcolorbox{black!50}{black!4}{\parbox{0.94\textwidth}{
\textbf{AI disclosure and review status.}
This audit used the OpenAI model Daybreak Blue at the ultra reasoning level.
OpenAI Codex performed AI-assisted target selection and attack-history
searching, scheme reconstruction, cryptanalysis, proof development,
implementation, experiment generation, and technical writing.  Separate
specialist AI agents performed adversarial review, and a fresh-context AI
agent independently reproduced the result with a separately written
implementation.  These were AI review and AI reproduction, not human
verification.  The technical, practical, adversarial AI,
independent-reproduction, final attack-history, and Publication Reviewer
gates all passed before public release.  No independent human verification
was performed or required.  OpenAI Codex generated and selected the title and
intentionally chose ``The Writing Was on the Floor'' as wordplay on
GLYPH/writing and on the zero $K$-floor that makes the forgery work.}}

\begin{abstract}
We give a public-key-only, zero-signing-query fresh-message forgery against
the final revision of Arjun Chopra's \emph{GLYPH: A New Instantiation of the
GLP Digital Signature Scheme}, IACR ePrint 2017/766.  The target claims
137.17 bits against its modeled forgery attack at the recommended row
$(m,n,q,B,\omega)=(2048,1024,59393,16383,16)$.  Its verifier hashes the
$K$-floor of $az_1+z_2-tc$, where $K=B-\omega=16367$, and checks that the
result encodes to the supplied sparse challenge $c$.

For any attacker-chosen fresh message, we first derive $c$ by hashing the
all-zero floor polynomial.  We set $z_1=0$ and, coefficient by coefficient,
choose $z_2$ from $\{0,-16366,16366\}$ so that every coefficient of
$z_2-tc$ modulo $q$ lies in the zero-floor interval $[0,32735)$.  The response
lies strictly inside the paper's $K$ bound.  The ordinary verifier therefore
hashes exactly the zero polynomial used to derive $c$ and accepts.  The
construction is deterministic once the public key and message are fixed,
uses no secret, and makes no signing query.

The sealed primary implementation accepted 64/64 fresh-key, fresh-message
forgeries under the pinned author's unmodified \texttt{glp\_verify}; all
512 negative controls rejected.  A separately written fresh-context AI
reproduction accepted another 64/64 full-parameter forgeries and rejected
all 256 controls.  Each reproduction exhaustively checked the selector over
all 59,393 residues.  The primary attack core took 68.648--206.126
microseconds, with a mean of 81.584 microseconds.  We located no previous
public substantive attack against final-v3 GLYPH in searches completed
through 13 October 2026.  The claim concerns the numerical signature and
message/hash interface specified by the paper and author source; neither
defines an author-standardized portable signature wire format.
\end{abstract}

\section{Introduction}

GLYPH is a 2017--2018 re-instantiation of the GLP lattice signature
scheme~\cite{glyph,glp}.  It retains GLP's Ring-LWE-style key relation and
Fiat--Shamir structure, changes the recommended parameters, and introduces
\textsc{Compress2} and a new $K$-floor operation.  The target presents its
recommended row as providing 137 bits of conservative security and estimates
the cost of its modeled forgery attack as $2^{137.17}$ basic classical
operations~\cite{glyph}.

The verifier's floor operation creates a much simpler route.  A forger need
not recover a secret or solve the short-vector problem modeled by the target.
It can decide in advance which floor polynomial the verifier will hash,
derive the challenge from that public value and a chosen message, and then
place each response coefficient into the corresponding floor bucket.  For
the recommended modulus and bounds, the zero bucket is reachable for every
possible coefficient of $-tc$ using one of three strictly bounded values.

The resulting attack has the following properties:
\begin{itemize}
  \item it targets the final ePrint object and the complete recommended
  parameter row, without weakening any parameter;
  \item it works for every public key and any attacker-chosen message;
  \item it produces the numerical tuple $(z_1,z_2,c)$ accepted by the
  ordinary verifier and uses the target's normal message-hash path;
  \item it requests no signature and uses no private value; and
  \item it is a design-level failure of the specified $K$-floor verification
  relation, rather than a parser, randomness, fault, or compatibility defect.
\end{itemize}

\paragraph{Technique attribution.}
We make no novelty claim for constructing a transcript backward from a
public verification equation, for choosing a public hash target before
solving for a response, or for manipulating a rounding bucket.  The target
itself analyzes a generic forgery strategy that chooses a hash target and
then seeks a short preimage with BKZ~\cite{glyph}.  Our contribution is the
observation and full-parameter demonstration that GLYPH's specific
$K$-floor relation makes such a preimage available coefficient by
coefficient without lattice reduction.

There is also earlier, distinct compression-related work in the GLP lineage.
Ducas, Durmus, Lepoint, and Lyubashevsky observed that the original GLP12
compression workaround permits changing a least-significant bit in an
already valid signature while preserving the checked high bits with high
probability~\cite{ddll}.  They explicitly identify this as a failure of
strong unforgeability.  That method starts with an existing signature and
produces another signature on the same message; it predates GLYPH's
\textsc{Compress2}/$K$-floor construction.  It is therefore not the
zero-query fresh-message attack proved here.

\section{Exact target, history, and claim boundary}

\subsection{Pinned final revision}

The target is \emph{GLYPH: A New Instantiation of the GLP Digital Signature
Scheme} by Arjun Chopra, IACR ePrint 2017/766~\cite{glyph}.  The audit uses
the current 14-page final object, operationally called final v3, last revised
7 June 2018.  Its exact PDF has SHA-256 digest
\begin{center}
\small\ttfamily
ff0eeaea3a9fcb830197ea954472869a08adbf42e64e727470fa8e0038c6682a.
\end{center}
The extracted text has SHA-256 digest
\begin{center}
\small\ttfamily
8ecd84282a5a752e1c2256597b9a10bdf50bfdb8eb2c572fa34eba27e54c1f2a.
\end{center}

The paper links the author's public-domain implementation.  We pin that
repository at commit
\begin{center}
\small\ttfamily
74b96d691391f687fe8f91e1d5cb3c251762d7cc
\end{center}
with Git tree
\begin{center}
\small\ttfamily
2ae8eec4b23dc57fc816c8935da8c2a59a753196.
\end{center}
The retained source archive has SHA-256 digest
\begin{center}
\small\ttfamily
43c68dfab7e717bf37ba57be9aa2fc49270f8eb42bfa420c85bf473aca12fe60.
\end{center}
The pinned source uses the same parameters and implements the paper's
numerical verification equation.

\subsection{Claimed security and recommended parameters}

Table~1 gives the sole recommended GLYPH row.  The paper instantiates the
hash with SHA-256, maps its output to a challenge of Hamming weight
$\omega=16$, and reports a 1.8~kB signature, 0.3~kB secret key, and 2.0~kB
public key.  It estimates primal key recovery, dual key recovery, and forgery
at $2^{139.5}$, $2^{175.8}$, and $2^{137.17}$ basic classical operations,
respectively.  These are the target's estimates, not claims revalidated by
this audit.  The attack below bypasses each modeled lattice problem.

\begin{table}[ht]
\centering
\caption{The audited full parameter row.}
\begin{tabular}{lrrrrrrrr}
\toprule
$m$ & $n$ & $q$ & $B$ & $\omega$ & $K=B-\omega$ & $D=2K+1$ & claimed forgery \\
\midrule
2048 & 1024 & 59393 & 16383 & 16 & 16367 & 32735 & $2^{137.17}$ \\
\bottomrule
\end{tabular}
\end{table}

\subsection{Dated attack-history result}

We located no previous public substantive attack against final-v3 GLYPH in
our searches completed through 13 October 2026.  The final search repeated
exact-title, ePrint-identifier, author, repository, $K$-floor, zero-query,
fresh-message, forgery, attack, and cryptanalysis queries; checked the
official record and author repository; screened the five citing works
reported by OpenAlex; and revisited the GLP lineage.  This is a bounded,
dated negative search result, not an absolute priority claim.

The closest located result is the original-GLP same-message malleability just
described~\cite{ddll}.  Located induced-fault work attacks original GLP under
faults rather than final-v3 GLYPH under ordinary execution~\cite{bindel}.
Masking work that used the GLYPH code base expressly focused on
countermeasures and did not report this numerical-verifier attack~\cite{barthe}.

\subsection{Intended interface}

The target paper defines a signature as the numerical tuple $(z_1,z_2,c)$ and
defines verification directly over that tuple and a message.  The author
source exposes the same tuple as an in-memory structure, accepts message
bytes, and implements the same hash and floor equation.  This is sufficient
to establish a design-level attack on the declared numerical signature
interface.

Neither the final paper nor the author source specifies a portable external
signature byte format or parser.  In particular, the source's
\texttt{poly2bytes} operation serializes a polynomial for hashing; it is not
an external signature codec.  We therefore make no claim about an unspecified
document wrapper, protocol envelope, or deployment-specific wire encoding.
The experimental codecs transport the numerical fields without changing
their semantics and are not presented as author-standardized formats.

\section{The GLYPH verifier}

\subsection{Ring and floor notation}

For the recommended power-of-two conductor, let
\[
  \R=\mathbb{Z}[x]/(x^{1024}+1),
  \qquad
  \Rq=\R/q\R,
  \qquad q=59393.
\]
Coefficients of elements in $\Rq$ are represented modulo $q$.  A centered
coefficient has its representative in
$[-(q-1)/2,(q-1)/2]$.

Set $K=B-\omega=16367$ and $D=2K+1=32735$.  For an integer $x$, write the
canonical residue $x\bmod q$ uniquely as
\[
  x\bmod q=rD+s,
  \qquad 0\le s\le 2K.
\]
GLYPH defines $\lfloor x\rfloor_K=r$ and applies this operation coefficient
by coefficient to ring elements.  Thus every residue in $[0,D)$ has
$K$-floor zero.

\subsection{Keys, signing, and verification}

The public domain parameter is $a\in\Rq$.  Key generation samples ternary
$s,e\in\Rq$ and publishes
\[
  t=as+e,
\]
while $(s,e)$ remains secret.  To sign a message $\mu$, the signer samples
bounded $y_1,y_2$, computes
\[
  c=\Encode\!\left(\Hash\!\left(
    \Kfloor{ay_1+y_2}\,\|\,\mu
  \right)\right),
\]
forms $z_1=sc+y_1$ and $z_2=ec+y_2$, rejects samples outside the coefficient
bound, and applies \textsc{Compress2} to $z_2$.  The challenge encoder
$\Encode$ returns a ring element with exactly $\omega=16$ coefficients in
$\{+1,-1\}$ and all remaining coefficients zero.

Given $(\mu,z_1,z_2,c)$, the verifier first checks the coefficient bounds for
$z_1$ and $z_2$.  It then computes
\begin{equation}
  h=az_1+z_2-tc\pmod q,
  \label{eq:reconstruct}
\end{equation}
and accepts exactly when
\begin{equation}
  c=\Encode\!\left(\Hash\!\left(
      \Kfloor{h}\,\|\,\mu
  \right)\right).
  \label{eq:verify}
\end{equation}
The author's source follows this path in \texttt{glp\_verify}, using two
little-endian bytes per floor coefficient as the internal hash encoding,
then hashing the message bytes and applying its sparse encoder.

Correctly generated signatures satisfy
\[
  az_1+z_2-tc=(as+e)c+ay_1+y_2-(as+e)c=ay_1+y_2
\]
before compression, and \textsc{Compress2} is intended to preserve the
$K$-floor.  Verification therefore reconstructs the same hash input used by
the signer.

\section{Zero-query fresh-message forgery}

\subsection{Attack algorithm}

Define
\[
  L=K-1=16366,
  \qquad D=2K+1=32735,
  \qquad q-L=43027.
\]
The public-key-only forger takes $(a,t)$ and an attacker-chosen message
$\mu$ as input:
\begin{enumerate}
  \item Let $w=0\in\Rq$ and compute
  \[
    c=\Encode\!\left(\Hash(w\,\|\,\mu)\right)
  \]
  using GLYPH's prescribed hash serialization and sparse encoding.
  \item Set $z_1=0\in\Rq$.
  \item Compute the public ring product $-tc\in\Rq$.  For each coefficient,
  let $u\in[0,q)$ be its canonical representative and set
  \begin{equation}
  z_2(u)=
  \begin{cases}
    0,  & 0\le u<D,\\
    -L, & D\le u<q-L,\\
    +L, & q-L\le u<q.
  \end{cases}
  \label{eq:selector}
  \end{equation}
  \item Output $(z_1,z_2,c)$ as the signature of $\mu$.
\end{enumerate}

The algorithm never invokes the signer.  Its only nontrivial algebra is
multiplication of the public polynomial $t$ by the public sparse challenge
$c$.

\subsection{Strict-interior selector proof}

\begin{lemma}
For every residue $u\in[0,q)$, the value $z_2(u)$ from
Equation~\eqref{eq:selector} satisfies $|z_2(u)|<K$ and
\[
  (u+z_2(u))\bmod q\in[0,D).
\]
\end{lemma}

\begin{proof}
The only possible centered response values are $0$ and $\pm L$, and
$L=K-1$, so the strict coefficient bound holds.  The three selector branches
give the following reconstructed intervals:
\[
\begin{array}{ccl}
0\le u<D
  &\Longrightarrow& u+0\in[0,32734],\\[2pt]
D\le u<q-L
  &\Longrightarrow& u-L\in[16369,26660],\\[2pt]
q-L\le u<q
  &\Longrightarrow& u+L-q\in[0,16365].
\end{array}
\]
Each interval lies inside $[0,D)=[0,32735)$.
\end{proof}

\begin{theorem}
For every GLYPH public key at the recommended row and every message $\mu$,
the public algorithm above produces a signature accepted by the ordinary
GLYPH verifier.  It makes zero signing queries and uses no secret input.
\end{theorem}

\begin{proof}
The response values meet even the paper's early strict $K$-bounded wording,
and $z_1=0$ is also strictly bounded.  The verifier's reconstructed
polynomial is
\[
  h=az_1+z_2-tc=z_2-tc.
\]
For each coefficient, $u=(-tc)_i\bmod q$.  The lemma shows that
$(u+(z_2)_i)\bmod q$ belongs to $[0,D)$, so its $K$-floor is zero.  Hence
\[
  \Kfloor{h}=0=w.
\]
The verifier's challenge recomputation in Equation~\eqref{eq:verify} is
therefore
\[
  \Encode\!\left(\Hash(0\,\|\,\mu)\right)=c,
\]
which is exactly how the attacker chose $c$.  The verifier accepts.
\end{proof}

\paragraph{No Fiat--Shamir circularity.}
The attacker does not search for a hash fixed point and does not program the
hash.  It chooses the verifier's intended floor output $w=0$ first, hashes
that public value with $\mu$ to obtain $c$, and only then solves the public
coefficient relation for $z_2$.  The selector succeeds for every possible
coefficient of the resulting $-tc$.

\paragraph{Freshness and success probability.}
The attacker chooses $\mu$ and never submits it, or any other message, to a
signing oracle.  It therefore produces a fresh-message signature with zero
queries.  The proof covers every public key, every challenge output, and
every message.  In the mathematical interface the success probability is
one; no favorable-key, favorable-message, or random-residue event is needed.

\subsection{Complexity}

The challenge $c$ has $\omega=16$ nonzero coefficients.  Computing $tc$ as
a sum of signed negacyclic shifts takes $O(n\omega)$ modular operations; the
selector and serialization take $O(n)$ work.  Hashing and sparse encoding
add their ordinary fixed-parameter costs.  Memory use is $O(n)$.  At
$n=1024$ this is a small public computation, in contrast with the target's
modeled block-$469$ BKZ forgery.

\section{Full-parameter experimental evidence}

Both reproductions use the complete recommended row and fresh OS-random
keys.  Together they cover 64 primary plus 64 independently reproduced
forgeries.  The measured experiments support the implementation claim; the
all-residue proof above establishes generality beyond the sampled keys.

\subsection{Primary reproduction}

The primary harness compiles the pinned author source without cryptographic
changes.  OpenSSL~3 no longer exports the low-level
\texttt{AES\_ctr128\_encrypt} entry point used by the 2018 source, so the
harness supplies a narrow API adapter through OpenSSL's
\texttt{CRYPTO\_ctr128\_encrypt} and \texttt{AES\_encrypt}.  Before any
attack trial, the upstream deterministic vector passes, exercising key
derivation, signing, verification, hashing, sparse encoding, compression,
and the adapter.

Key generation, forging, and verification run as separate processes.  The
isolated forger accepts only a canonical public key and a message.  Its
binary links neither the GLYPH key-generation/signing object nor any signing
symbol.  The separately linked verifier contains the pinned
\texttt{glp\_verify} symbol and no signing symbol.  No secret is serialized
or passed to the forger, and the recorded signing-query count is zero.

For each of 64 trials, the driver creates a fresh OS-random key and a fresh
attacker message, invokes the public forger, and submits the tuple to the
ordinary verifier.  All 64 forgeries accept.  All $64\cdot1024=65{,}536$
reconstructed coefficients have zero floor.  Public-key, message, signature,
and challenge fingerprints are distinct in all 64 trials.

\begin{table}[ht]
\centering
\caption{Sealed full-parameter results.}
\begin{tabular}{lrr}
\toprule
Measurement & Primary & Fresh-context reproduction \\
\midrule
fresh OS-random keys/messages & 64 & 64 \\
ordinary-verifier accepts & 64/64 & 64/64 \\
zero-floor coefficients & 65,536/65,536 & 65,536/65,536 \\
negative-control rejects & 512/512 & 256/256 \\
selector residues checked & 59,393 & 59,393 \\
selector failures & 0 & 0 \\
signing queries & 0 & 0 \\
workers & 1 & 1 \\
\bottomrule
\end{tabular}
\end{table}

The primary attack-core time ranged from 68,648 to 206,126 nanoseconds; the
median was 76,994 nanoseconds and the mean was 81,584 nanoseconds.  The
complete 64-trial driver took 1.16 seconds wall-clock and used 53,760~KiB
peak resident memory.  It used one worker and exited successfully.

\subsection{Primary controls}

Each trial applies eight negative controls.  Changed, truncated, and extended
messages all reject.  Flipping a challenge sign and changing a challenge
position reject.  A response change that remains within the coefficient
bound but forces a nonzero floor rejects.  Finally, out-of-range $z_1$ and
$z_2$ coefficients reject at the verifier's bound check.  These controls
exercise the hash binding, sparse-challenge comparison, floor relation, and
response bounds; they rule out an accidentally permissive harness.

The driver also enumerates every $u\in[0,59393)$ for the strict-interior
selector in Equation~\eqref{eq:selector}.  It separately checks the author's
inclusive \textsc{Compress2} boundary selector.  Both exhaustive checks pass
all 59,393 residues with zero failure.  The attack uses only the
strict-interior selector.

\subsection{Fresh-context independent AI reproduction}

A fresh-context AI agent received the public target pins, public attack
description, and reproduction requirements.  It did not read the primary
package.  It independently implemented challenge derivation, negacyclic
sparse multiplication, and coefficient selection in C, then used the pinned
author implementation only for legitimate setup and the ordinary
\texttt{glp\_verify} call.

For 64 fresh OS-random keys and 48-byte fresh messages, all 64 independent
forgeries accepted.  Changed-message, changed-challenge, in-bound
changed-floor, and out-of-bound controls rejected in every trial, for 256/256
rejections.  Its independently written selector also passed all 59,393
residues, and all 65,536 reconstructed trial coefficients had zero floor.
The experiment body took 2.596779495 seconds; external timing was 3.50
seconds with 5,120~KiB peak resident memory and one worker.

The combined independent setup/reproduction binary necessarily links the
complete upstream object, so unused signing symbols are present.  Source and
disassembly review found no call to \texttt{glp\_sign} or
\texttt{glp\_deterministic\_sign}.  The independent forging function accepts
only a public key and message, and the harness cleanses the setup secret
before calling it.  This is fresh-context AI reproduction, not independent
human verification.

\subsection{Evidence integrity}

The primary machine-readable results, package manifest, and checksum list
have SHA-256 digests
\begin{center}
\begin{tabular}{ll}
\toprule
Artifact & SHA-256 \\
\midrule
primary \texttt{results.json}
& \ttfamily\footnotesize 3d48ec90fdab886323e189359f45e46851113eb9d74a36b76dccc5219a2acc52 \\
primary \texttt{manifest.json}
& \ttfamily\footnotesize 37f35f099485fd61cdd001a3ac76425d06c141bf5e3b9a87951ade100cabf329 \\
primary \texttt{MANIFEST.sha256}
& \ttfamily\footnotesize 595e948beecd42200e38a36b0d6c370b474c55fa5de90ca593b53d023816f53d \\
independent \texttt{MANIFEST.sha256}
& \ttfamily\footnotesize f2ce4a22e576f662bfa1b761c806916c60d120bd6dd58fddd5691f716621c6a7 \\
independent source
& \ttfamily\footnotesize 4fb82963255a7e4ac50c742c29591163e6a899481adf6964d84b91218a60b37b \\
\bottomrule
\end{tabular}
\end{center}
All 37 entries in the primary checksum list and all 52 entries in the
independent list verify.  The primary isolated-forger source has SHA-256
digest
\begin{center}
\small\ttfamily
26d69afea2af2ea7f82fe56dbf77460419eb906dabc1b965928bb53f96fde095.
\end{center}

\section{Why the verification equation fails}

GLYPH's modeled forgery in Section~5.3 of the target begins with a uniform
ring element $y$, hashes $(y,\mu)$ to obtain $c$, and asks BKZ to find bounded
$z_1,z_2$ satisfying $az_1+z_2=tc+y$.  The security estimate treats matching
the complete ring element $y$ as a short-preimage problem.

The actual verifier does not hash $az_1+z_2-tc$ itself.  It hashes only its
$K$-floor.  Every one of the 59,393 possible residues can be shifted into the
zero-floor bucket by a response in $\{0,-(K-1),K-1\}$.  Setting $z_1=0$
therefore decomposes the apparent ring preimage problem into 1,024 independent
three-choice operations.  Challenge sparsity makes the remaining public
product especially cheap, but the proof does not rely on favorable
challenge positions or signs.

The weakness is in the mathematical verification relation.  The OpenSSL
adapter only restores an API removed since publication and is covered by the
upstream vector.  The attack uses valid numerical fields and ordinary hash
and challenge derivation; it does not exploit malformed input, a parser,
nonce reuse, weak randomness, a fault, timing leakage, or undefined behavior.

\section{Repair discussion and limitations}

Preventing $z_1=0$ alone does not repair the underlying issue: the security
argument must account for the quotient information revealed by $K$-floor and
show that an attacker cannot choose a floor target first and solve for a
bounded response afterward.  Likewise, banning the three response values
used here would address one instantiation of the selector rather than the
public reachability problem.

A repair must redesign the transcript binding or compression relation so
that producing a bounded response for a challenge derived from an
attacker-selected verifier image is computationally hard.  Any proposed
change should be accompanied by a proof for the exact compressed verifier,
including modular wraparound and every accepted response value, and by a
fresh parameter analysis.  Adding an independently specified outer document
hash does not change this attack: the construction works for whatever message
bytes reach the declared GLYPH hash interface.

Our result is limited in the following ways:
\begin{itemize}
  \item it attacks the final-v3 numerical tuple and message/hash interface
  specified by the pinned paper and source;
  \item it does not claim an author-standardized external signature wire
  format, because none is defined by those artifacts;
  \item it does not claim secret-key recovery, signing-key recovery, or an
  attack on every GLP-family design; and
  \item it makes no novelty claim for verification-first construction,
  public rounding-bucket manipulation, or the earlier GLP12 same-message
  malleability method.
\end{itemize}

The exact result is stronger than a same-message alternate-signature attack:
the adversary selects a fresh message and needs no prior signature.  It is
also independent of the paper's ambiguity between early strict
$K$-bounded wording and the later inclusive verifier boundary, because all
constructed response coefficients have magnitude at most $K-1$.

\section{Reproducibility}

The public package provides two entry points.  From the attack directory,
\texttt{python3 reproduce.py} performs the quick eight-trial full-parameter
run, while \texttt{python3 reproduce-all-rows.py} performs the complete
64-trial run.  Each entry point rebuilds the pinned upstream verifier and the
public C forger, exhaustively checks all 59,393 selector residues, creates
fresh keys and messages, submits every forgery to the ordinary verifier, runs
all controls, records resources and per-trial public digests, and verifies its
checksum manifest.  The underlying commands are
\texttt{./reproducer/reproduce.sh --quick} and
\texttt{./reproducer/reproduce.sh --full}.

Reproduction requires a C compiler, Python~3 for the public wrapper,
OpenSSL development headers and library, POSIX shell utilities, and GNU
\texttt{time} for the reported external resource record.  Both modes use
one worker.  Random keys and messages intentionally change across runs, so
the expected invariant is the reported acceptance, rejection, bound, floor,
and exhaustive-residue counts rather than byte-identical trial fingerprints.
The target PDF, source commit and tree, source-file hashes, commands,
environment, binaries, symbol reports, results, and manifests are retained by
each run so that the exact evidence boundary can be checked independently.

\section{Conclusion}

At GLYPH's complete recommended parameter row, every coefficient of the
public product $-tc$ can be moved into the verifier's zero $K$-floor bucket
with one of three strictly bounded response values.  A forger can consequently
choose any fresh
message, derive the public sparse challenge from an all-zero floor target,
set $z_1=0$, and solve for $z_2$ coefficient by coefficient.  The ordinary
verifier accepts for every key and message, with zero signing queries.

The sealed primary and fresh-context AI reproductions accepted 64+64
full-parameter forgeries, rejected 512+256 controls, and independently
exhausted all 59,393 selector residues.  The primary attack core averaged
81.584 microseconds, bypassing the target's claimed $2^{137.17}$-operation
forgery cost.  The result is a practical design-level fresh-message forgery
against final-v3 GLYPH's declared numerical signature interface.

\section*{Acknowledgements}
The computations were performed on the Norwegian Research and Education Cloud (NREC), using resources provided by the University of Bergen and the University of Oslo.

\begin{thebibliography}{9}

\bibitem{glyph}
A.~Chopra.
\newblock GLYPH: A New Instantiation of the GLP Digital Signature Scheme.
\newblock IACR Cryptology ePrint Archive, Report 2017/766, final revision,
2018.
\newblock \url{https://eprint.iacr.org/2017/766}.

\bibitem{glp}
T.~G\"uneysu, V.~Lyubashevsky, and T.~P\"oppelmann.
\newblock Practical Lattice-Based Cryptography: A Signature Scheme for
Embedded Systems.
\newblock In \emph{Cryptographic Hardware and Embedded Systems---CHES 2012},
LNCS 7428, pp.~530--547. Springer, 2012.
\newblock \url{https://doi.org/10.1007/978-3-642-33027-8_31}.

\bibitem{ddll}
L.~Ducas, A.~Durmus, T.~Lepoint, and V.~Lyubashevsky.
\newblock Lattice Signatures and Bimodal Gaussians.
\newblock In \emph{CRYPTO 2013, Part I}, LNCS 8042, pp.~40--56. Springer,
2013; IACR ePrint Report 2013/383.
\newblock \url{https://eprint.iacr.org/2013/383}.

\bibitem{bindel}
N.~Bindel, J.~Buchmann, and J.~Kr\"amer.
\newblock Lattice-Based Signature Schemes and Their Sensitivity to Fault
Attacks.
\newblock IACR Cryptology ePrint Archive, Report 2016/415, 2016.
\newblock \url{https://eprint.iacr.org/2016/415}.

\bibitem{barthe}
G.~Barthe, S.~Bela\"id, T.~Espitau, P.-A.~Fouque, B.~Gr\'egoire,
M.~Rossi, and M.~Tibouchi.
\newblock Masking the GLP Lattice-Based Signature Scheme at Any Order.
\newblock In \emph{EUROCRYPT 2018, Part II}, LNCS 10821, pp.~354--384.
Springer, 2018; IACR ePrint Report 2018/381.
\newblock \url{https://eprint.iacr.org/2018/381}.

\end{thebibliography}

\end{document}
